% Insert after Online Expert Data Collection, before Policy Update.
% This specifies the directional-kernel, bounded-wait variant tested overnight.
% It does not attribute fixed-gate results to this adaptive method.
\subsection{Expert-Feedback Timing Calibration}
\label{sec:ef_timing_calibration}

The ID-calibrated PCA gate provides an initial intervention rule. We refine
its timing using the expert behavior observed after each takeover. Our target
is a region where assistance begins to provide behavior that differs from
the current policy, while requiring little reversal of preceding robot
motion. Rather than assigning a fixed preferred timestep to a task, we use
completed interventions to infer whether subsequent takeovers near a visited
state should occur earlier, remain unchanged, or be postponed locally.
The policy parameters, PCA subspace, and baseline threshold $\gamma_B$
remain fixed throughout the collection round.

\paragraph{Action disagreement and corrective motion.}
For intervention $i$, let $t_i$ denote its onset and
$\{(o_{i,k},a^*_{i,k})\}_{k=0}^{L_i-1}$ the expert suffix, with each
observation recorded before its corresponding action. At these observations,
we freely generate action chunks from the frozen policy and compare the
first $\Delta$ actions, corresponding to one execution block:
\begin{equation}
e_{i,k}=\frac{1}{M\Delta d_a}
\sum_{m=1}^{M}\sum_{h=0}^{\Delta-1}
\|\widehat a^{(m)}_{\theta,h}(o_{i,k})-a^*_{i,k+h}\|_2^2.
\label{eq:ef_disagreement}
\end{equation}
Here $d_a$ is the number of active action dimensions, and predictions and
expert actions use the same execution scaling. Predictions are generated
without conditioning on the expert action chunk. We calibrate $q_e$ as the
95th percentile of trajectory-maximum disagreement on ID demonstrations,
using the same inference protocol. The first crossing
$k_i^*=\min\{k:e_{i,k}>q_e\}$ localizes the onset of elevated disagreement
along the observed expert path. Complete target windows are required for
this feedback; an unobserved crossing is treated as censored.

We also measure how much the first expert execution block reverses the
preceding policy block. Let $v_i^-,v_i^+$ be their actual end-effector
displacements and $b_i^-,b_i^+$ their changes in gripper width. Define
\begin{equation}
R_i=\sqrt{
 \left(\frac{[-\langle v_i^+,v_i^-\rangle/\|v_i^-\|_2]_+}{\sigma_p}\right)^2
 +\left(\frac{[-\operatorname{sgn}(b_i^-)b_i^+]_+}{\sigma_w}\right)^2}.
\label{eq:ef_reversal}
\end{equation}
A component is zero when its preceding motion has zero magnitude; the cue
is unavailable without complete preceding and following blocks. Its reference
$q_R$ is calibrated from adjacent blocks of successful ID demonstrations.
Disagreement and reversal serve different roles: the former describes when
the policy departs from observed expert behavior, while the latter provides
evidence of initial corrective motion.

\paragraph{Local timing feedback.}
If $R_i>q_R$, we attach an earlier-takeover cue $y_i=+1$ to the preceding
observed policy-query state. Otherwise, an initial low-disagreement prefix
covering at least one execution block provides a postponement cue $y_i=-1$
at the actual takeover state. This includes censored suffixes only when
the required complete low-disagreement windows have been observed. When
disagreement appears within the first block without excessive reversal, we
store $y_i=0$ at the takeover state: assistance is already locally warranted,
but the observation does not justify moving it earlier. Corrective-motion
cues take precedence. Insufficient observations produce no cue.

Each cue stores the original bridge feature $z_i^B$ of its attributed policy
state. This implements a local preference for the operational region
\begin{equation}
\mathcal W_\theta(\tau)
=\{t:k^*(t)<\Delta,\ R(t)\le q_R\},
\label{eq:ef_preferred_region}
\end{equation}
where $k^*(t)$ and $R(t)$ describe the expert continuation from $t$.
The algorithm observes only actual interventions, and transfers their
feedback through feature similarity. It neither enumerates all such
continuations nor assumes a unique preferred time.

\paragraph{Regularized threshold adjustment.}
At a subsequent query, let $w_{ti}$ be a Gaussian similarity weight between
ID-normalized bridge features, truncated outside a calibrated radius. With
one cue per intervention, we fit a local timing direction by
\begin{equation}
\widehat d_t
=\arg\min_d\left\{\lambda d^2+
\sum_{i\in\mathcal M}w_{ti}(d-y_i)^2\right\}
=\frac{\sum_i w_{ti}y_i}{\lambda+\sum_i w_{ti}}.
\label{eq:ef_local_direction}
\end{equation}
The zero-centered prior retains the baseline under weak support. We set
$\widehat d_t=0$ when fewer than two interventions support the query or
the absolute weighted mean direction is below $0.25$. The threshold is
\begin{equation}
\gamma_t=\gamma_B\exp(-\beta\widehat d_t).
\label{eq:ef_local_threshold}
\end{equation}
Positive feedback lowers the threshold and negative feedback raises it;
because $|\widehat d_t|\le1$, the adjustment is bounded. The gate evaluates
$s_{\mathrm{BPCA}}(o_t)>\gamma_t$ before executing each action block.
If feedback suppresses a baseline alarm, it authorizes at most one additional
block; assistance is requested at the next query unless the episode has
already terminated. This prevents repeated reuse of one local wait cue from
causing indefinite postponement.

Feedback is computed after the expert segment and committed before the next
episode. Empty memory recovers the original PCA gate, and subsequent policy
updates follow the same data-retention and fine-tuning procedure. No extra
expert timing labels are requested. The preferred region describes our
feedback objective; its relationship to learning efficiency is evaluated
through downstream policy improvement under a specified expert budget.
